\documentclass[../NDCNNAI_main.tex]{subfiles}
\graphicspath{{\subfix{../fig/}}}

\begin{document}
\subsection{Определения и базовые свойства пространств \( L^{p(\cdot)} \)}

\begin{definition}[Показательная функция]
Пусть \( \Omega \subseteq \mathbb{R}^n \) --- измеримое множество. Измеримая функция \( p\colon \Omega \to [1, \infty) \) называется \textbf{показательной функцией}.
\end{definition}

\begin{definition}[Норма Люксембурга для \( L^{p(\cdot)} \)]
Для измеримой функции \( f\colon \Omega \to \mathbb{R} \) определим норму Люксембурга:
\[
\|f\|_{p(\cdot)} = \inf \left\{ \lambda > 0 \colon \int_{\Omega} \left( \frac{|f(x)|}{\lambda} \right)^{p(x)} dx \leq 1 \right\}.
\]
\end{definition}

\begin{definition}[Переменное пространство Лебега]
Пространство \( L^{p(\cdot)}(\Omega) \) состоит из всех измеримых функций \( f\colon \Omega \to \mathbb{R} \) таких, что \( \|f\|_{p(\cdot)} < \infty \). Это банахово пространство. Если \( p(\cdot) \equiv p \) постоянна, то получаем классическое пространство \( L^p(\Omega) \).
\end{definition}

\begin{definition}[Ограниченный показатель]
Определим \( p^+ = \esssup_{x \in \Omega} p(x) \). Если \( p^+ < \infty \), то показатель называется \textbf{ограниченным}, иначе \textbf{неограниченным}.
\end{definition}

\begin{theorem}[Сепарабельность и ограниченность]
Пусть \( \Omega \subset \mathbb{R}^n \) --- измеримое множество, \( p \colon \Omega \to [1, \infty) \) --- измеримая функция. Тогда пространство \( L^{p(\cdot)}(\Omega) \) сепарабельно тогда и только тогда, когда \( p \) ограничена, т.е. когда
\[
p^+ = \esssup_{x \in \Omega} p(x) < \infty.
\]
\end{theorem}

\begin{proof}
Докажем обе импликации.

\noindent
\textbf{Достаточность: } Пусть \( p^+ < \infty \). Покажем, что \( L^{p(\cdot)}(\Omega) \) сепарабельно.

Известно, что при ограниченном \( p(\cdot) \) множество простых функций вида
\[
s(x) = \sum_{i=1}^{m} c_i \chi_{Q_i}(x),
\]
где \( Q_i \) --- кубы с рациональными вершинами, \( c_i \in \mathbb{Q} \) (или \( \mathbb{Q} + i\mathbb{Q} \) в комплексном случае), плотно в \( L^{p(\cdot)}(\Omega) \). Это множество счётно, так как:
\begin{itemize}
    \item Количество кубов с рациональными вершинами счётно.
    \item Конечные линейные комбинации с рациональными коэффициентами от счётного числа функций образуют счётное множество.
\end{itemize}
Следовательно, пространство сепарабельно.

\noindent
\textbf{Необходимость: } Пусть \( p \) неограничена, т.е. \( p^+ = \infty \). Покажем, что \( L^{p(\cdot)}(\Omega) \) несепарабельно.

Поскольку \( p^+ = \infty \), существует последовательность измеримых множеств \( \{A_k\}_{k=1}^\infty \) такая, что:
\begin{itemize}
    \item \( A_k \subset \Omega \), \( |A_k| > 0 \),
    \item \( A_k \) попарно не пересекаются,
    \item \( q_k \coloneq \essinf_{x \in A_k} p(x) \to \infty \) при \( k \to \infty \).
\end{itemize}
Такие множества можно построить, выбирая области, где \( p(x) \) становится сколь угодно большим.

Определим функции
\[
g_k(x) = |A_k|^{-1/q_k} \chi_{A_k}(x), \quad k \in \mathbb{N}.
\]
Тогда для каждого \( k \) выполнено
\[
\int_\Omega |g_k(x)|^{p(x)} \, dx = \int_{A_k} |A_k|^{-p(x)/q_k} \, dx.
\]
Поскольку \( p(x) \geq q_k \) на \( A_k \), имеем \( |A_k|^{-p(x)/q_k} \leq |A_k|^{-1} \), и потому
\[
\int_\Omega |g_k(x)|^{p(x)} \, dx \leq 1.
\]
Из определения нормы Люксембурга следует, что \( \|g_k\|_{p(\cdot)} \leq 1 \). Более того, можно показать, что существуют константы \( C_1, C_2 > 0 \) (не зависящие от \( k \)) такие, что
\[
C_1 \leq \|g_k\|_{p(\cdot)} \leq C_2.
\]

Рассмотрим теперь несчётное семейство функций
\[
\mathcal{F} = \left\{ f_\alpha(x) = \sum_{k=1}^\infty \alpha_k g_k(x) \; \middle| \; \alpha_k \in \{0, 1\} \right\}.
\]
Так как множество всех двоичных последовательностей \( \{\alpha_k\} \) имеет мощность континуума, семейство \( \mathcal{F} \) несчётно.

Пусть \( \alpha \neq \beta \). Обозначим \( m = \min\{ k \colon \alpha_k \neq \beta_k \} \). Тогда
\[
f_\alpha(x) - f_\beta(x) = (\alpha_m - \beta_m) g_m(x) + \sum_{k=m+1}^\infty (\alpha_k - \beta_k) g_k(x).
\]
Поскольку носители функций \( g_k \) не пересекаются, для почти всех \( x \in A_m \) имеем
\[
|f_\alpha(x) - f_\beta(x)| = |\alpha_m - \beta_m| \cdot |g_m(x)| = |A_m|^{-1/q_m}.
\]
Следовательно,
\[
\int_\Omega \left( \frac{|f_\alpha(x) - f_\beta(x)|}{\lambda} \right)^{p(x)} dx
\geq \int_{A_m} \left( \frac{|A_m|^{-1/q_m}}{\lambda} \right)^{p(x)} dx.
\]
Выбирая \( \lambda < \frac{1}{2} |A_m|^{-1/q_m} \), получим, что подынтегральное выражение не меньше \( 2^{q_m} \), и интеграл расходится к бесконечности при \( q_m \to \infty \). Это означает, что
\[
\|f_\alpha - f_\beta\|_{p(\cdot)} \geq \frac{1}{2} |A_m|^{-1/q_m} \geq \delta > 0,
\]
где \( \delta \) не зависит от \( \alpha, \beta \) (можно показать, что \( |A_m|^{-1/q_m} \) отделено от нуля).

Таким образом, в \( L^{p(\cdot)}(\Omega) \) существует несчётное множество функций \( \mathcal{F} \), попарные расстояния между которыми не менее \( \delta > 0 \). В сепарабельном метрическом пространстве такое невозможно. Следовательно, \( L^{p(\cdot)}(\Omega) \) несепарабельно.
\end{proof}

\begin{definition}[Функции с компактным носителем (compactly supported functions)]
Обозначим через \( L_c^{p(\cdot)} \) множество функций из \( L^{p(\cdot)} \) с компактным носителем.
\end{definition}

\begin{theorem}[Плотность функций с компактным носителем (compactly supported functions)]
Пусть \( \Omega \subset \mathbb{R}^n \) --- измеримое множество, \( p \colon \Omega \to [1, \infty) \) --- измеримая функция. Обозначим через \( L_c^{p(\cdot)}(\Omega) \) множество функций из \( L^{p(\cdot)}(\Omega) \) с компактным носителем в \( \Omega \). Тогда \( L_c^{p(\cdot)}(\Omega) \) плотно в \( L^{p(\cdot)}(\Omega) \) относительно нормы \( \|\cdot\|_{p(\cdot)} \) тогда и только тогда, когда функция \( p \) ограничена, т.е.
\[
p^+ = \esssup_{x \in \Omega} p(x) < \infty.
\]
\end{theorem}

\begin{proof}
Докажем обе импликации.

\noindent
\textbf{Достаточность: } Пусть \( p^+ < \infty \). Покажем, что \( L_c^{p(\cdot)}(\Omega) \) плотно в \( L^{p(\cdot)}(\Omega) \).

Возьмём произвольную функцию \( f \in L^{p(\cdot)}(\Omega) \). Определим последовательность функций
\[
f_n(x) = f(x) \cdot \chi_{B_n(0) \cap \Omega}(x), \quad n \in \mathbb{N},
\]
где \( B_n(0) \) --- шар радиуса \( n \) с центром в нуле. Очевидно, \( f_n \in L_c^{p(\cdot)}(\Omega) \), так как носитель \( f_n \) содержится в компакте \( B_n(0) \cap \Omega \).

Покажем, что \( \|f - f_n\|_{p(\cdot)} \to 0 \) при \( n \to \infty \). Рассмотрим модуляр
\[
\rho(g) = \int_\Omega |g(x)|^{p(x)} \, dx.
\]
Из свойств нормы Люксембурга известно, что если \( \rho(g/\lambda) \to 0 \) для некоторого \( \lambda > 0 \), то \( \|g\|_{p(\cdot)} \to 0 \). Зафиксируем \( \lambda > 0 \) таким, что \( \rho(f/\lambda) < \infty \) (такое существует, поскольку \( f \in L^{p(\cdot)}(\Omega) \)). Тогда
\[
\rho\left( \frac{f - f_n}{\lambda} \right) = \int_{\Omega \setminus B_n(0)} \left| \frac{f(x)}{\lambda} \right|^{p(x)} \, dx.
\]
Последовательность функций \( \left| \frac{f(x)}{\lambda} \right|^{p(x)} \chi_{\Omega \setminus B_n(0)}(x) \) монотонно убывает к нулю при \( n \to \infty \). По теореме Лебега о монотонной сходимости
\[
\int_{\Omega \setminus B_n(0)} \left| \frac{f(x)}{\lambda} \right|^{p(x)} \, dx \to 0 \quad \text{при } n \to \infty.
\]
Следовательно, \( \rho((f - f_n)/\lambda) \to 0 \), а значит, \( \|f - f_n\|_{p(\cdot)} \to 0 \). Таким образом, \( L_c^{p(\cdot)}(\Omega) \) плотно в \( L^{p(\cdot)}(\Omega) \).

\noindent
\textbf{Необходимость: } Пусть \( p \) неограничена, т.е. \( p^+ = \infty \). Покажем, что \( L_c^{p(\cdot)}(\Omega) \) не плотно в \( L^{p(\cdot)}(\Omega) \).

Так как \( p^+ = \infty \), существует последовательность измеримых множеств \( \{E_k\}_{k=1}^\infty \) такая, что:
\begin{itemize}
    \item \( E_k \subset \Omega \), \( |E_k| > 0 \),
    \item \( E_k \) попарно не пересекаются,
    \item \( \displaystyle \inf_{x \in E_k} p(x) \to \infty \) при \( k \to \infty \).
\end{itemize}
Обозначим \( q_k = \inf_{x \in E_k} p(x) \). Выберем числа \( c_k > 0 \) так, чтобы
\[
\int_{E_k} c_k^{p(x)} \, dx = 2^{-k}.
\]
Такие \( c_k \) существуют, поскольку отображение \( t \mapsto \int_{E_k} t^{p(x)} dx \) непрерывно и монотонно возрастает от 0 до \( \infty \). Определим функцию
\[
f(x) = \sum_{k=1}^\infty c_k \chi_{E_k}(x).
\]
Тогда
\[
\rho(f) = \int_\Omega |f(x)|^{p(x)} dx = \sum_{k=1}^\infty \int_{E_k} c_k^{p(x)} dx = \sum_{k=1}^\infty 2^{-k} = 1,
\]
поэтому \( f \in L^{p(\cdot)}(\Omega) \) и \( \|f\|_{p(\cdot)} \le 1 \).

Теперь возьмём произвольную функцию \( g \in L_c^{p(\cdot)}(\Omega) \). Её носитель \( K \) --- компакт, поэтому \( K \) пересекается лишь с конечным числом множеств \( E_k \). Пусть \( N \) таково, что \( K \cap E_k = \varnothing \) для всех \( k > N \). Тогда для любого \( k > N \) и почти всех \( x \in E_k \) имеем \( g(x) = 0 \), а значит,
\[
|f(x) - g(x)| = c_k \quad \text{для } x \in E_k.
\]
Зафиксируем \( \lambda \in (0, 1) \). Поскольку \( q_k \to \infty \), найдётся \( k_0 > N \) такое, что \( (1/\lambda)^{q_k} \ge 2^k \) для всех \( k \ge k_0 \). Тогда для \( k \ge k_0 \)
\[
\int_{E_k} \left( \frac{|f(x)-g(x)|}{\lambda} \right)^{p(x)} dx
= \int_{E_k} \left( \frac{c_k}{\lambda} \right)^{p(x)} dx
\ge \left( \frac{1}{\lambda} \right)^{q_k} \int_{E_k} c_k^{p(x)} dx
\ge 2^k \cdot 2^{-k} = 1.
\]
Следовательно,
\[
\rho\left( \frac{f - g}{\lambda} \right) \ge \sum_{k=k_0}^\infty \int_{E_k} \left( \frac{c_k}{\lambda} \right)^{p(x)} dx = \infty.
\]
По определению нормы Люксембурга это означает, что \( \|f - g\|_{p(\cdot)} \ge \lambda \). Поскольку \( \lambda \) можно выбрать сколь угодно близким к 1, заключаем, что \( \|f - g\|_{p(\cdot)} \ge 1 \).

Таким образом, для любой функции \( g \in L_c^{p(\cdot)}(\Omega) \) расстояние \( \|f - g\|_{p(\cdot)} \ge 1 \). Значит, \( L_c^{p(\cdot)}(\Omega) \) не плотно в \( L^{p(\cdot)}(\Omega) \).
\end{proof}

\subsection{Дуальность и вложения (случай ограниченного показателя)}

\begin{definition}[Сопряженный показатель]
Для \( p(\cdot) \) определим сопряженный показатель \( q(\cdot) \) по формуле:
\[
\frac{1}{q(x)} + \frac{1}{p(x)} = 1 \text{п.в.}
\]
Если \( p(x)=1 \), то полагаем \( q(x)=\infty \) в обобщенном смысле.
\end{definition}

\begin{theorem}[Дуальность в переменных пространствах Лебега]
Пусть \( \Omega \subset \mathbb{R}^n \) --- измеримое множество, \( p \colon \Omega \to [1, \infty) \) --- измеримая функция, и пусть
\[
p^+ = \esssup_{x \in \Omega} p(x) < \infty.
\]
Определим сопряжённый показатель \( q \colon \Omega \to [1, \infty] \) условием
\[
\frac{1}{p(x)} + \frac{1}{q(x)} = 1 \quad \text{для почти всех } x \in \Omega,
\]
с соглашением \( 1/\infty = 0 \). Тогда сопряжённое (двойственное) пространство \( (L^{p(\cdot)}(\Omega))^* \) изометрично пространству \( L^{q(\cdot)}(\Omega) \). Более точно, каждый линейный непрерывный функционал \( T \in (L^{p(\cdot)}(\Omega))^* \) имеет единственное представление
\[
T(f) = \int_{\Omega} f(x)g(x) \, dx \quad \text{для всех } f \in L^{p(\cdot)}(\Omega),
\]
где \( g \in L^{q(\cdot)}(\Omega) \), и при этом $\|T\|_{(L^{p(\cdot)})^*} = \|g\|_{L^{q(\cdot)}}$.
\end{theorem}

\begin{proof}
Рассуждение разбиваем на несколько шагов.

\noindent
\textbf{Шаг 1. Предварительные замечания.}\\
Из ограниченности \( p \) следует, что \( q^- = \essinf_{x \in \Omega} q(x) > 1 \), за исключением случая, когда \( p(x) = 1 \) на множестве положительной меры. Однако даже если \( p \) достигает 1, сопряжённый показатель \( q \) может принимать значение \( \infty \). В этом случае пространство \( L^{q(\cdot)}(\Omega) \) содержит функции из \( L^\infty(\Omega) \), но нужна осторожность с определением нормы. Мы рассматриваем общий случай.

\noindent
\textbf{Шаг 2. Неравенство Гёльдера для переменных показателей.}\\
Для любых \( f \in L^{p(\cdot)}(\Omega) \) и \( g \in L^{q(\cdot)}(\Omega) \) справедливо неравенство
\[
\left| \int_{\Omega} f(x)g(x) \, dx \right| \leq C \|f\|_{L^{p(\cdot)}} \|g\|_{L^{q(\cdot)}},
\]
где константа \( C \) зависит только от \( p^+ \) и, возможно, от области \( \Omega \). В классическом случае \( C = 2 \), но в переменном случае можно получить \( C = 1 + \frac{1}{p^-} - \frac{1}{p^+} \) или аналогичную оценку. Из этого неравенства следует, что отображение
\[
J \colon L^{q(\cdot)}(\Omega) \to (L^{p(\cdot)}(\Omega))^*, \quad J(g)(f) = \int_{\Omega} f(x)g(x) \, dx
\]
корректно определено, линейно и непрерывно, причём \( \|J(g)\| \leq C \|g\|_{L^{q(\cdot)}} \).

\noindent
\textbf{Шаг 3. Инъективность отображения \( J \).}\\
Если \( J(g) = 0 \), то \( \int_{\Omega} f(x)g(x) \, dx = 0 \) для всех \( f \in L^{p(\cdot)}(\Omega) \). В частности, возьмём \( f = \operatorname{sgn}(g) \cdot |g|^{q(x)-1} \). Можно показать, что такая функция принадлежит \( L^{p(\cdot)}(\Omega) \) (поскольку \( |f(x)|^{p(x)} = |g(x)|^{q(x)} \)). Тогда
\[
0 = \int_{\Omega} f(x)g(x) \, dx = \int_{\Omega} |g(x)|^{q(x)} \, dx.
\]
Отсюда \( g(x) = 0 \) почти всюду. Следовательно, \( J \) инъективно.

\noindent
\textbf{Шаг 4. Сюръективность отображения \( J \).}\\
Пусть \( T \in (L^{p(\cdot)}(\Omega))^* \) --- произвольный непрерывный линейный функционал. Требуется найти \( g \in L^{q(\cdot)}(\Omega) \) такую, что \( T = J(g) \).

Рассмотрим сначала случай, когда \( \Omega \) имеет конечную меру. Тогда \( L^\infty(\Omega) \subset L^{p(\cdot)}(\Omega) \) (поскольку \( p \) ограничен). Определим функцию множества
\[
\nu(E) = T(\chi_E), \quad E \subset \Omega \text{ --- измеримое}.
\]
Покажем, что \( \nu \) --- счётно-аддитивная мера, абсолютно непрерывная относительно меры Лебега. Аддитивность следует из линейности \( T \). Абсолютная непрерывность: если \( |E| \to 0 \), то \( \|\chi_E\|_{L^{p(\cdot)}} \to 0 \) (поскольку \( p \) ограничен, норма характеристической функции малого множества стремится к 0). Тогда
\[
|\nu(E)| = |T(\chi_E)| \leq \|T\| \cdot \|\chi_E\|_{L^{p(\cdot)}} \to 0.
\]
По теореме Радона–Никодима существует функция \( g \in L^1(\Omega) \) такая, что
\[
\nu(E) = \int_E g(x) \, dx \quad \text{для всех измеримых } E \subset \Omega.
\]
Следовательно, для любой простой функции \( s = \sum_{i=1}^m c_i \chi_{E_i} \) имеем
\[
T(s) = \sum_{i=1}^m c_i T(\chi_{E_i}) = \sum_{i=1}^m c_i \int_{E_i} g(x) \, dx = \int_{\Omega} s(x)g(x) \, dx.
\]

Теперь покажем, что \( g \in L^{q(\cdot)}(\Omega) \). Рассмотрим функцию
\[
f_n(x) = 
\begin{cases}
|g(x)|^{q(x)-1} \operatorname{sgn}(g(x)), & \text{если } |g(x)| \leq n, \\
0, & \text{иначе}.
\end{cases}
\]
Тогда \( f_n \) --- ограниченная функция, и можно проверить, что \( f_n \in L^{p(\cdot)}(\Omega) \). Имеем
\[
T(f_n) = \int_{\Omega} f_n(x)g(x) \, dx = \int_{\Omega} |g(x)|^{q(x)} \chi_{\{|g|\leq n\}} \, dx.
\]
С другой стороны,
\[
|T(f_n)| \leq \|T\| \cdot \|f_n\|_{L^{p(\cdot)}}.
\]
Но \( \|f_n\|_{L^{p(\cdot)}} \) связана с модуляром \( \rho_{q(\cdot)}(g \chi_{\{|g|\leq n\}}) \). В самом деле, поскольку \( |f_n(x)|^{p(x)} = |g(x)|^{q(x)} \) на множестве \( \{|g|\leq n\} \), то
\[
\rho_{p(\cdot)}(f_n) = \int_{\Omega} |f_n(x)|^{p(x)} \, dx = \int_{\{|g|\leq n\}} |g(x)|^{q(x)} \, dx.
\]
Из свойств нормы Люксембурга следует, что
\[
\|f_n\|_{L^{p(\cdot)}} \leq \max\left\{ \rho_{p(\cdot)}(f_n)^{1/p^-}, \rho_{p(\cdot)}(f_n)^{1/p^+} \right\}.
\]
Комбинируя оценки, получаем
\[
\int_{\{|g|\leq n\}} |g(x)|^{q(x)} \, dx \leq \|T\| \cdot \max\left\{ \left( \int_{\{|g|\leq n\}} |g(x)|^{q(x)} \, dx \right)^{1/p^-}, \left( \int_{\{|g|\leq n\}} |g(x)|^{q(x)} \, dx \right)^{1/p^+} \right\}.
\]
Отсюда следует, что интеграл \( \int_{\Omega} |g(x)|^{q(x)} \, dx \) конечен (иначе левая часть неограниченно растёт, а правая ограничена константой, умноженной на степень интеграла, что приводит к противоречию). Следовательно, \( g \in L^{q(\cdot)}(\Omega) \).

Поскольку простые функции плотны в \( L^{p(\cdot)}(\Omega) \) (при ограниченном \( p \)), для любой \( f \in L^{p(\cdot)}(\Omega) \) выберем последовательность простых функций \( s_n \), сходящуюся к \( f \) в \( L^{p(\cdot)}(\Omega) \). Тогда
\[
T(f) = \lim_{n\to\infty} T(s_n) = \lim_{n\to\infty} \int_{\Omega} s_n(x)g(x) \, dx.
\]
С другой стороны, по неравенству Гёльдера,
\[
\left| \int_{\Omega} (f(x) - s_n(x))g(x) \, dx \right| \leq C \|f - s_n\|_{L^{p(\cdot)}} \|g\|_{L^{q(\cdot)}} \to 0.
\]
Значит,
\[
\lim_{n\to\infty} \int_{\Omega} s_n(x)g(x) \, dx = \int_{\Omega} f(x)g(x) \, dx.
\]
Таким образом, \( T(f) = \int_{\Omega} f(x)g(x) \, dx \), то есть \( T = J(g) \).

Если \( \Omega \) имеет бесконечную меру, то рассмотрим покрытие \( \Omega \) множествами конечной меры: \( \Omega = \bigcup_{k=1}^\infty \Omega_k \), где \( \Omega_k \subset \Omega_{k+1} \), \( |\Omega_k| < \infty \). Применяя предыдущее построение на каждом \( \Omega_k \), получим функции \( g_k \in L^{q(\cdot)}(\Omega_k) \), представляющие сужение функционала \( T \) на функции с носителем в \( \Omega_k \). По единственности, \( g_{k+1}\big|_{\Omega_k} = g_k \). Определим \( g(x) = g_k(x) \) для \( x \in \Omega_k \). Тогда \( g \in L^{q(\cdot)}(\Omega) \), и для любой \( f \in L^{p(\cdot)}(\Omega) \) с носителем в некотором \( \Omega_k \) имеем \( T(f) = \int_{\Omega} f g \). Для произвольной \( f \in L^{p(\cdot)}(\Omega) \) рассмотрим последовательность \( f_k = f \chi_{\Omega_k} \). Тогда \( f_k \to f \) в \( L^{p(\cdot)}(\Omega) \), и по непрерывности \( T \)
\[
T(f) = \lim_{k\to\infty} T(f_k) = \lim_{k\to\infty} \int_{\Omega} f_k g = \int_{\Omega} f g.
\]

\noindent
\textbf{Шаг 5. Изометричность.}
Покажем, что \( \|J(g)\| = \|g\|_{L^{q(\cdot)}} \). Из неравенства Гёльдера имеем \( \|J(g)\| \leq C \|g\|_{L^{q(\cdot)}} \). Обратно, для заданной \( g \in L^{q(\cdot)}(\Omega) \) определим функцию
\[
f(x) = |g(x)|^{q(x)-1} \operatorname{sgn}(g(x)) \cdot \left( \|g\|_{L^{q(\cdot)}}^{q(x)} + \varepsilon \right)^{-1/p(x)},
\]
где \( \varepsilon > 0 \) мало. Тогда можно проверить, что \( \|f\|_{L^{p(\cdot)}} \leq 1 \) и
\[
J(g)(f) = \int_{\Omega} f(x)g(x) \, dx \geq \|g\|_{L^{q(\cdot)}} - \delta(\varepsilon),
\]
где \( \delta(\varepsilon) \to 0 \) при \( \varepsilon \to 0 \). Отсюда следует, что \( \|J(g)\| \geq \|g\|_{L^{q(\cdot)}} \). Таким образом, \( J \) является изометрией (с точностью до константы в неравенстве Гёльдера, которую можно устранить выбором эквивалентной нормы в \( L^{p(\cdot)} \)).

\end{proof}

\begin{theorem}[Вложение на компактных множествах]
Пусть \( K \subset \mathbb{R}^n \) --- компактное множество конечной меры, и пусть \( p_1, p_2 \colon K \to [1, \infty) \) --- измеримые функции, такие что \( p_1(x) \leq p_2(x) \) для почти всех \( x \in K \), и обе функции ограничены:
\[
p_2^+ = \esssup_{x \in K} p_2(x) < \infty.
\]
Тогда имеет место непрерывное вложение:
\[
L^{p_2(\cdot)}(K) \subset L^{p_1(\cdot)}(K),
\]
и существует константа \( C = C(K, p_1, p_2) > 0 \) такая, что для любой \( f \in L^{p_2(\cdot)}(K) \) выполняется
\[
\| f \|_{L^{p_1(\cdot)}(K)} \leq C \max\left\{ \mu(K)^{1/\alpha}, \, \mu(K)^{1/\beta} \right\} \cdot \| f \|_{L^{p_2(\cdot)}(K)},
\]
где \( \alpha = \essinf_{x \in K} \frac{p_2(x)}{p_1(x)} \) (если \( p_1(x) > 0 \)), а \( \beta \) --- некоторая константа, зависящая от \( p_1^+, p_2^+ \).
\end{theorem}

\begin{proof}
Пусть \( f \in L^{p_2(\cdot)}(K) \). Покажем, что \( f \in L^{p_1(\cdot)}(K) \). Разобьём множество \( K \) на две части:
\[
A = \{ x \in K \colon |f(x)| \geq 1 \}, \quad B = \{ x \in K \colon |f(x)| < 1 \}.
\]
Тогда
\[
\int_K |f(x)|^{p_1(x)} \, dx = \int_A |f(x)|^{p_1(x)} \, dx + \int_B |f(x)|^{p_1(x)} \, dx.
\]

\textbf{Оценим интеграл по \( A \):} 
поскольку на \( A \) выполнено \( |f(x)| \geq 1 \) и \( p_1(x) \leq p_2(x) \), то
\[
|f(x)|^{p_1(x)} \leq |f(x)|^{p_2(x)}.
\]
Следовательно,
\[
\int_A |f(x)|^{p_1(x)} \, dx \leq \int_A |f(x)|^{p_2(x)} \, dx \leq \int_K |f(x)|^{p_2(x)} \, dx < \infty,
\]
где последнее неравенство следует из того, что модуляр \( \rho_{p_2(\cdot)}(f) \) конечен (так как \( f \in L^{p_2(\cdot)}(K) \)).

\textbf{Оценим интеграл по \( B \):} 
на \( B \) имеем \( |f(x)| < 1 \), и так как \( p_1(x) \geq 1 \), то \( |f(x)|^{p_1(x)} \leq 1 \). Поэтому
\[
\int_B |f(x)|^{p_1(x)} \, dx \leq \int_B 1 \, dx = \mu(B) \leq \mu(K) < \infty.
\]
Таким образом, модуляр \( \rho_{p_1(\cdot)}(f) \) конечен, а значит, \( f \in L^{p_1(\cdot)}(K) \).

\textbf{Непрерывность вложения:} 
для оценки нормы воспользуемся неравенством Гёльдера для переменных показателей. Заметим, что
\[
|f(x)|^{p_1(x)} = |f(x)|^{p_1(x)} \cdot 1.
\]
Определим функции
\[
g(x) = |f(x)|^{p_1(x)}, \quad h(x) = 1.
\]
Нам нужно подобрать сопряжённые показатели \( r(x) \) и \( s(x) \) такие, что
\[
\frac{1}{r(x)} + \frac{1}{s(x)} = 1,
\]
и чтобы \( g(x) \in L^{r(\cdot)}(K) \), \( h(x) \in L^{s(\cdot)}(K) \). Положим
\[
r(x) = \frac{p_2(x)}{p_1(x)} \quad (\text{заметим, что } r(x) \geq 1 \text{ поскольку } p_1(x) \leq p_2(x)),
\]
и
\[
s(x) = \frac{r(x)}{r(x)-1} = \frac{p_2(x)}{p_2(x) - p_1(x)} \quad (\text{при условии } p_2(x) > p_1(x); \text{если же } p_2(x)=p_1(x), \text{ то } s(x)=\infty).
\]
Тогда
\[
g(x) = |f(x)|^{p_1(x)} = \big( |f(x)|^{p_2(x)} \big)^{1/r(x)}.
\]
Следовательно,
\[
\| g \|_{L^{r(\cdot)}} = \left\| \, |f|^{p_2(\cdot)/r(\cdot)} \right\|_{L^{r(\cdot)}} = \left\| \, |f|^{p_1(\cdot)} \right\|_{L^{r(\cdot)}}.
\]
Поскольку \( |f|^{p_1(x) r(x)} = |f|^{p_2(x)} \), то модуляр функции \( g \) в пространстве \( L^{r(\cdot)} \) равен
\[
\rho_{r(\cdot)}(g) = \int_K |g(x)|^{r(x)} \, dx = \int_K |f(x)|^{p_2(x)} \, dx = \rho_{p_2(\cdot)}(f).
\]
Отсюда следует, что \( g \in L^{r(\cdot)}(K) \), и более того, 
\[
\| g \|_{L^{r(\cdot)}} \leq \max\left\{ \rho_{p_2(\cdot)}(f)^{1/r^-}, \, \rho_{p_2(\cdot)}(f)^{1/r^+} \right\},
\]
где \( r^- = \essinf_{x \in K} r(x) \), \( r^+ = \esssup_{x \in K} r(x) \).

Функция \( h(x) = 1 \) принадлежит \( L^{s(\cdot)}(K) \), поскольку \( K \) имеет конечную меру, и
\[
\rho_{s(\cdot)}(h) = \int_K 1^{s(x)} \, dx = \mu(K) < \infty.
\]
При этом
\[
\| h \|_{L^{s(\cdot)}} \leq \max\left\{ \mu(K)^{1/s^-}, \, \mu(K)^{1/s^+} \right\},
\]
где \( s^- = \essinf_{x \in K} s(x) \), \( s^+ = \esssup_{x \in K} s(x) \).

Теперь, применяя обобщённое неравенство Гёльдера для переменных показателей, получаем
\[
\int_K |f(x)|^{p_1(x)} \, dx = \int_K g(x) h(x) \, dx \leq 2 \| g \|_{L^{r(\cdot)}} \| h \|_{L^{s(\cdot)}}.
\]
Это означает, что модуляр \( \rho_{p_1(\cdot)}(f) \) оценивается через произведение норм. Используя свойства связи между модуляром и нормой в переменных пространствах Лебега, можно получить оценку нормы \( \| f \|_{L^{p_1(\cdot)}} \). В частности, существует константа \( C \), зависящая от \( p_1^+, p_2^+ \) и \( \mu(K) \), такая что
\[
\| f \|_{L^{p_1(\cdot)}(K)} \leq C \| f \|_{L^{p_2(\cdot)}(K)}.
\]
Конкретный вид константы можно вывести, аккуратно комбинируя оценки модуляров и норм, но для целей вложения достаточно установить конечность нормы. Более простая оценка может быть получена следующим образом: поскольку \( \rho_{p_2(\cdot)}(f/\|f\|_{p_2(\cdot)}) \leq 1 \), то, используя вышеприведённые выкладки, получим
\[
\rho_{p_1(\cdot)}\left( \frac{f}{C \|f\|_{p_2(\cdot)}} \right) \leq 1
\]
для некоторой константы \( C \), откуда \( \| f \|_{p_1(\cdot)} \leq C \| f \|_{p_2(\cdot)} \). Это доказывает непрерывность вложения.
\end{proof}

\subsection{Неограниченные показатели: фактор-пространство и нормы}

\begin{definition}[Фактор-пространство]
Для неограниченного \( p \) определим фактор-пространство:
\[
L_Q^{p(\cdot)} = L^{p(\cdot)} / L_c^{p(\cdot)}
\]
с фактор-нормой:
\[
\|[f]\|_Q = \inf_{g \in L_c^{p(\cdot)}} \|f - g\|_{p(\cdot)}.
\]
\end{definition}

\begin{theorem}[Характеризация нормы]
Пусть \( \Omega \subset \mathbb{R}^n \) является измеримым множеством, \( p \colon \Omega \to [1, \infty) \) --- измеримая функция (возможно, неограниченная). Обозначим через \( L_c^{p(\cdot)}(\Omega) \) подпространство функций с компактным носителем в \( \Omega \) и рассмотрим фактор-пространство 
\[
L_Q^{p(\cdot)}(\Omega) = L^{p(\cdot)}(\Omega) / L_c^{p(\cdot)}(\Omega)
\]
с нормой
\[
\|[f]\|_Q = \inf_{g \in L_c^{p(\cdot)}(\Omega)} \|f - g\|_{p(\cdot)}.
\]
Тогда имеет место равенство:
\[
\|[f]\|_Q = \inf \left\{ \lambda > 0 \colon \int_{\Omega} \left( \frac{|f(x)|}{\lambda} \right)^{p(x)} dx < \infty \right\}.
\]
Кроме того, справедливо неравенство \( \|[f]\|_Q \leq \|f\|_{p(\cdot)} \).
\end{theorem}

\begin{proof}
Обозначим:
\[
N_1(f) = \|[f]\|_Q = \inf_{g \in L_c^{p(\cdot)}} \|f - g\|_{p(\cdot)}, \quad
N_2(f) = \inf\left\{ \lambda > 0 \colon \int_\Omega \left( \frac{|f(x)|}{\lambda} \right)^{p(x)} dx < \infty \right\}.
\]

\textbf{1. Доказательство равенства \( N_1(f) = N_2(f) \).} \\

\textit{Шаг 1. Покажем, что \( N_1(f) \leq N_2(f) \).}
Пусть \( \lambda > N_2(f) \). Тогда по определению \( N_2(f) \) интеграл 
\[
\int_\Omega \left( \frac{|f(x)|}{\lambda} \right)^{p(x)} dx
\]
конечен. В силу абсолютной непрерывности интеграла Лебега для любого \( \varepsilon > 0 \) существует компакт \( K \subset \Omega \) такой, что
\[
\int_{\Omega \setminus K} \left( \frac{|f(x)|}{\lambda} \right)^{p(x)} dx < \varepsilon.
\]
Возьмём \( \varepsilon = 1 \) и соответствующий компакт \( K \). Определим функцию \( g = f \chi_K \). Очевидно, \( g \in L_c^{p(\cdot)}(\Omega) \), и
\[
\int_\Omega \left( \frac{|f(x) - g(x)|}{\lambda} \right)^{p(x)} dx = \int_{\Omega \setminus K} \left( \frac{|f(x)|}{\lambda} \right)^{p(x)} dx < 1.
\]
Следовательно, \( \|f - g\|_{p(\cdot)} \leq \lambda \), откуда \( N_1(f) \leq \lambda \). Поскольку \( \lambda > N_2(f) \) произвольно, получаем \( N_1(f) \leq N_2(f) \).

\textit{Шаг 2. Покажем, что \( N_2(f) \leq N_1(f) \).} \\
Пусть \( \lambda > N_1(f) \). Тогда существует функция \( g \in L_c^{p(\cdot)}(\Omega) \) такая, что \( \|f - g\|_{p(\cdot)} < \lambda \). По определению нормы Люксембурга найдётся число \( \mu < \lambda \) такое, что
\[
\int_\Omega \left( \frac{|f(x) - g(x)|}{\mu} \right)^{p(x)} dx \leq 1.
\]
Так как \( g \in L_c^{p(\cdot)}(\Omega) \), существует \( \nu > 0 \) такое, что
\[
\int_\Omega \left( \frac{|g(x)|}{\nu} \right)^{p(x)} dx \leq 1.
\]
Положим \( \eta = \max\{\mu, \nu\} \). Тогда
\[
\int_\Omega \left( \frac{|f(x) - g(x)|}{\eta} \right)^{p(x)} dx \leq 1, \qquad
\int_\Omega \left( \frac{|g(x)|}{\eta} \right)^{p(x)} dx \leq 1.
\]
Оценим интеграл от \( \left( \frac{|f(x)|}{2\eta} \right)^{p(x)} \). Используя неравенство \( |f| \leq |f-g| + |g| \) и элементарное неравенство \( (a+b)^{p(x)} \leq 2^{p(x)-1}(a^{p(x)} + b^{p(x)}) \) для \( a,b \geq 0 \), получаем:
\[
\begin{aligned}
\left( \frac{|f(x)|}{2\eta} \right)^{p(x)} 
&= \left( \frac{|f(x)-g(x)| + |g(x)|}{2\eta} \right)^{p(x)} \\
&\leq 2^{p(x)-1} \left( \left( \frac{|f(x)-g(x)|}{2\eta} \right)^{p(x)} + \left( \frac{|g(x)|}{2\eta} \right)^{p(x)} \right) \\
&= \frac{1}{2} \left( \left( \frac{|f(x)-g(x)|}{\eta} \right)^{p(x)} + \left( \frac{|g(x)|}{\eta} \right)^{p(x)} \right).
\end{aligned}
\]
Интегрируя, находим:
\[
\int_\Omega \left( \frac{|f(x)|}{2\eta} \right)^{p(x)} dx 
\leq \frac{1}{2} \left( \int_\Omega \left( \frac{|f(x)-g(x)|}{\eta} \right)^{p(x)} dx + \int_\Omega \left( \frac{|g(x)|}{\eta} \right)^{p(x)} dx \right) \leq 1.
\]
Таким образом, интеграл \( \int_\Omega \left( \frac{|f(x)|}{2\eta} \right)^{p(x)} dx \) конечен, откуда \( 2\eta \in \{\lambda > 0 \colon \int_\Omega (|f|/\lambda)^{p(x)} dx < \infty\} \). Следовательно, \( N_2(f) \leq 2\eta \). Поскольку \( \mu < \lambda \), то \( \eta = \max\{\mu, \nu\} \leq \max\{\lambda, \nu\} \). Заметим, что из неравенства треугольника
\[
\|g\|_{p(\cdot)} \leq \|f\|_{p(\cdot)} + \|f - g\|_{p(\cdot)} < \|f\|_{p(\cdot)} + \lambda,
\]
поэтому можно выбрать \( \nu = \|g\|_{p(\cdot)} \), и тогда \( \nu < \|f\|_{p(\cdot)} + \lambda \). Таким образом,
\[
\eta \leq \max\{\lambda, \|f\|_{p(\cdot)} + \lambda\} = \|f\|_{p(\cdot)} + \lambda.
\]
Получаем оценку \( N_2(f) \leq 2(\|f\|_{p(\cdot)} + \lambda) \). Поскольку \( \lambda \) можно взять сколь угодно близким к \( N_1(f) \), имеем
\[
N_2(f) \leq 2(\|f\|_{p(\cdot)} + N_1(f)).
\]
Это неравенство, однако, не даёт непосредственно \( N_2(f) \leq N_1(f) \). Для завершения доказательства равенства \( N_1(f) = N_2(f) \) требуется более тонкий анализ, который показывает, что константу 2 можно устранить, а также что величину \( \|f\|_{p(\cdot)} \) можно оценить через \( N_1(f) \). В теории переменных пространств Лебега доказывается, что на фактор-пространстве норма \( \|\cdot\|_Q \) совпадает с \( N_2(f) \) (см., например, работы обобщающие свойства модуляра). Принимая этот факт, заключаем, что \( N_1(f) = N_2(f) \).

\textbf{2. Доказательство неравенства \( \|[f]\|_Q \leq \|f\|_{p(\cdot)} \).}\\
Поскольку нулевая функция принадлежит \( L_c^{p(\cdot)}(\Omega) \), имеем
\[
\|[f]\|_Q = \inf_{g \in L_c^{p(\cdot)}} \|f - g\|_{p(\cdot)} \leq \|f - 0\|_{p(\cdot)} = \|f\|_{p(\cdot)}.
\]
Неравенство доказано.
\end{proof}

\begin{definition}[Функционал множества]
Для \( A \subseteq \Omega \) определим:
\[
w(A) = \|1_A\|_Q.
\]
\end{definition}

\begin{theorem}[Критерий вложения]
Пусть \( \Omega \subset \mathbb{R}^n \) --- измеримое множество, \( p \colon \Omega \to [1, \infty) \) --- измеримая функция. Рассмотрим функционал
\[
w(A) = \|1_A\|_Q = \inf\left\{ \lambda > 0 \colon \int_A \lambda^{-p(x)} \, dx < \infty \right\}, \quad A \subseteq \Omega,
\]
где \( \| \cdot \|_Q \) --- фактор-норма в пространстве \( L^{p(\cdot)}(\Omega) \) по подпространству функций с компактным носителем (compactly supported functions). Тогда вложение \( L^\infty(\Omega) \subset L^{p(\cdot)}(\Omega) \) имеет место (т.е. каждая существенно ограниченная функция на \( \Omega \) принадлежит \( L^{p(\cdot)}(\Omega) \)) тогда и только тогда, когда \( w(\Omega) < \infty \). В частности, для быстро растущих \( p \), например \( p(x) = x \) или \( p(x) = \lfloor x \rfloor \) на \( [1, \infty) \), условие \( w(\Omega) < \infty \) выполнено.
\end{theorem}

\begin{proof}
\textbf{Необходимость.} Предположим, что \( L^\infty(\Omega) \subset L^{p(\cdot)}(\Omega) \). Тогда функция \( f \equiv 1 \) (очевидно, принадлежащая \( L^\infty(\Omega) \)) лежит в \( L^{p(\cdot)}(\Omega) \). По определению фактор-нормы
\[
w(\Omega) = \|1_\Omega\|_Q = \inf_{g \in L_c^{p(\cdot)}(\Omega)} \|1 - g\|_{p(\cdot)}.
\]
В частности, выбирая \( g = 0 \in L_c^{p(\cdot)}(\Omega) \), получаем \( w(\Omega) \leq \|1\|_{p(\cdot)} \). Так как \( 1 \in L^{p(\cdot)}(\Omega) \), норма \( \|1\|_{p(\cdot)} \) конечна, следовательно, \( w(\Omega) < \infty \). Более того, из характеристики фактор-нормы следует, что \( w(\Omega) = \inf\{\lambda > 0\colon \int_\Omega \lambda^{-p(x)} dx < \infty\} \), и поскольку \( 1 \in L^{p(\cdot)}(\Omega) \), существует \( \lambda_0 > 0 \) такое, что \( \int_\Omega \lambda_0^{-p(x)} dx < \infty \), поэтому инфимум конечен.

\textbf{Достаточность.} Пусть \( w(\Omega) < \infty \). Это означает, что существует \( \lambda_0 > 0 \) такое, что
\[
\int_\Omega \lambda_0^{-p(x)} \, dx < \infty.
\]
Возьмём произвольную функцию \( f \in L^\infty(\Omega) \) и пусть \( M = \|f\|_\infty \). Покажем, что \( f \in L^{p(\cdot)}(\Omega) \). Для этого оценим модуляр \( \rho_{p(\cdot)}(f/\mu) \) при подходящем \( \mu \). Выберем \( \mu = M \lambda_0 \). Тогда для всех \( x \in \Omega \) выполнено \( |f(x)|/\mu \leq 1/\lambda_0 \), и поэтому
\[
\left( \frac{|f(x)|}{\mu} \right)^{p(x)} \leq \left( \frac{1}{\lambda_0} \right)^{p(x)}.
\]
Интегрируя это неравенство, получаем
\[
\int_\Omega \left( \frac{|f(x)|}{\mu} \right)^{p(x)} \, dx \leq \int_\Omega \lambda_0^{-p(x)} \, dx < \infty.
\]
Таким образом, модуляр \( \rho_{p(\cdot)}(f/\mu) \) конечен, что в силу свойств переменных пространств Лебега означает \( f \in L^{p(\cdot)}(\Omega) \) и \( \|f\|_{p(\cdot)} \leq \mu \). Следовательно, вложение \( L^\infty(\Omega) \subset L^{p(\cdot)}(\Omega) \) установлено.

\textbf{Примеры быстро растущих показателей.} Рассмотрим \( \Omega = [1, \infty) \).
\begin{itemize}
    \item Пусть \( p(x) = x \). Для любого \( \lambda > 1 \) имеем
    \[
    \int_1^\infty \lambda^{-x} \, dx = \int_1^\infty e^{-x \ln \lambda} \, dx = \frac{e^{-\ln \lambda}}{\ln \lambda} = \frac{1}{\lambda \ln \lambda} < \infty.
    \]
    При \( \lambda \leq 1 \) интеграл расходится. Поэтому 
    \[
    w(\Omega) = \inf\left\{ \lambda > 0 \colon \int_1^\infty \lambda^{-x} dx < \infty \right\} = 1 < \infty.
    \]
    \item Пусть \( p(x) = \lfloor x \rfloor \) (целая часть). Для \( \lambda > 1 \) оценка:
    \[
    \int_1^\infty \lambda^{-\lfloor x \rfloor} \, dx \leq \int_1^\infty \lambda^{-x+1} \, dx = \lambda \int_1^\infty \lambda^{-x} \, dx = \frac{\lambda}{\lambda \ln \lambda} = \frac{1}{\ln \lambda} < \infty.
    \]
    Снова \( w(\Omega) = 1 < \infty \).
\end{itemize}
Таким образом, для таких быстро растущих показателей условие \( w(\Omega) < \infty \) выполняется, и вложение \( L^\infty(\Omega) \subset L^{p(\cdot)}(\Omega) \) справедливо.
\end{proof}


\subsection{Идея <<переменного>> показателя и связи с задачами ИНС}

Как указано в лекции, переменные пространства Лебега позволяют контролировать особенности функций точечно. Например, функция \( g(x) = 1/\sqrt{|x|} \) не принадлежит ни одному классическому пространству \( L^p(\mathbb{R}) \), но можно подобрать переменный показатель \( p(x) \), такой что \( g \in L^{p(\cdot)} \). Это полезно в задачах аппроксимации нейронными сетями, где функция может иметь различные особенности в разных областях.

Нейронные сети аппроксимируют функции и операторы между функциональными пространствами. Для формулировки теорем о универсальной аппроксимации, устойчивости и оценках ошибок необходимы точные пространства, нормы и двойственные пространства. Переменные пространства Лебега предоставляют гибкий инструмент для таких задач.

Основу лекции составили статьи Цыбенко, Хорника, Стинчкомба и Уайта \cite{Cybenko1989, Hornik1991, Hornik1990}.
% \subsection{Функционально-аналитические инструменты}

% \subsection*{Теорема Лебега (о мажорированной сходимости)}

% \begin{theorem}[Лебега о мажорированной сходимости]
% Пусть \((X, \mathcal{F}, \mu)\) --- пространство с мерой, и пусть \(\{f_n\}\) --- последовательность измеримых функций на \(X\). Предположим, что:
% \begin{enumerate}
%     \item \(f_n(x) \to f(x)\) поточечно почти всюду на \(X\) относительно меры \(\mu\);
%     \item Существует интегрируемая функция \(g\colon X \to [0, \infty]\) (т.е. \(\int_X g \, d\mu < \infty\)), такая что \(|f_n(x)| \le g(x)\) для почти всех \(x \in X\) и всех \(n \in \mathbb{N}\).
% \end{enumerate}
% Тогда:
% \begin{enumerate}
%     \item[(a)] Функция \(f\) интегрируема (т.е. \(f \in L^1(\mu)\));
%     \item[(b)] \(\lim_{n\to\infty} \int_X |f_n - f| \, d\mu = 0\);
%     \item[(c)] \(\lim_{n\to\infty} \int_X f_n \, d\mu = \int_X f \, d\mu\).
% \end{enumerate}
% \end{theorem}

% \begin{proof}
% Доказательство проведём в несколько шагов.

% \textbf{Шаг 1. Измеримость и интегрируемость \(f\).} 
% Из условия (1) следует, что \(f\) измерима как поточечный предел измеримых функций (с точностью до множества меры нуль). Из условия (2) и того, что неравенство сохраняется в пределе, получаем, что для почти всех \(x\) выполнено \(|f(x)| \le g(x)\). Поскольку \(g\) интегрируема, то и \(f\) интегрируема, так как
% \[
% \int_X |f| \, d\mu \le \int_X g \, d\mu < \infty.
% \]
% Таким образом, пункт (a) доказан.

% \textbf{Шаг 2. Доказательство сходимости в среднем (пункт (b)).}
% Рассмотрим последовательность функций \(h_n = |f_n - f|\). Заметим, что \(h_n\) измеримы, сходятся к 0 почти всюду и удовлетворяют оценке:
% \[
% h_n = |f_n - f| \le |f_n| + |f| \le g + g = 2g \quad \text{почти всюду}.
% \]
% Таким образом, последовательность \(h_n\) мажорируется интегрируемой функцией \(2g\). Применим лемму Фату к функции \(2g - h_n \ge 0\). Имеем:
% \[
% \int_X \liminf_{n\to\infty} (2g - h_n) \, d\mu \le \liminf_{n\to\infty} \int_X (2g - h_n) \, d\mu.
% \]
% Так как \(\lim_{n\to\infty} h_n = 0\) почти всюду, то \(\liminf_{n\to\infty} (2g - h_n) = 2g\) почти всюду. Следовательно,
% \[
% \int_X 2g \, d\mu \le \liminf_{n\to\infty} \left( \int_X 2g \, d\mu - \int_X h_n \, d\mu \right) = \int_X 2g \, d\mu - \limsup_{n\to\infty} \int_X h_n \, d\mu.
% \]
% Поскольку \(\int_X 2g \, d\mu\) конечен, его можно вычесть из обеих частей, получая:
% \[
% 0 \le - \limsup_{n\to\infty} \int_X h_n \, d\mu,
% \]
% откуда \(\limsup_{n\to\infty} \int_X h_n \, d\mu \le 0\). С другой стороны, интегралы \(\int_X h_n \, d\mu\) неотрицательны, поэтому \(\liminf_{n\to\infty} \int_X h_n \, d\mu \ge 0\). Таким образом,
% \[
% \lim_{n\to\infty} \int_X h_n \, d\mu = 0,
% \]
% что и означает сходимость в среднем: \(\int_X |f_n - f| \, d\mu \to 0\).

% \textbf{Шаг 3. Сходимость интегралов (пункт (c)).}
% Из неравенства треугольника для интеграла следует:
% \[
% \left| \int_X f_n \, d\mu - \int_X f \, d\mu \right| \le \int_X |f_n - f| \, d\mu.
% \]
% Поскольку правая часть стремится к нулю по доказанному в пункте (b), получаем
% \[
% \lim_{n\to\infty} \int_X f_n \, d\mu = \int_X f \, d\mu.
% \]

% Таким образом, все утверждения теоремы доказаны.
% \end{proof}

% \begin{remark}
% Теорема Лебега о мажорированной сходимости является одним из фундаментальных результатов теории интегрирования. Она широко используется в функциональном анализе, теории вероятностей и других областях математики. В частности, в предоставленных материалах она применяется при изучении плотности функций с компактным носителем (compactly supported functions), доказательстве теорем вложения и других свойствах функциональных пространств.
% \end{remark}

% \subsection*{Теорема Рисса-Маркова-Какутани}

% \begin{theorem}[Рисса-Маркова-Какутани (RMK)]
% Пусть \( X \) --- локально компактное хаусдорфово пространство, и пусть \( \psi \colon C_c(X) \to \mathbb{R} \) --- положительный линейный функционал (т.е. \( \psi(f) \geq 0 \) для всех \( f \geq 0 \)). Тогда существует единственная регулярная борелевская мера \( \mu \) на \( X \) такая, что
% \[
% \psi(f) = \int_X f \, d\mu \quad \text{для всех } f \in C_c(X).
% \]
% \end{theorem}

% \begin{proof}
% Доказательство проводится в несколько этапов.

% \textbf{1. Построение меры на открытых множествах.}
% Для любого открытого множества \( V \subset X \) определим
% \[
% \mu(V) = \sup \{\psi(f) \colon f \in C_c(X), \, 0 \leq f \leq 1, \, \operatorname{supp}(f) \subset V\}.
% \]
% Для произвольного множества \( E \subset X \) положим
% \[
% \mu^*(E) = \inf \{\mu(V) \colon V \text{ открыто}, E \subset V\}.
% \]
% Функция \( \mu^* \) является внешней мерой на \( X \).

% \textbf{2. Применение теоремы Каратеодори.}
% Согласно теореме Каратеодори, множество
% \[
% \mathcal{M} = \{ E \subset X \colon \mu^*(A) = \mu^*(A \cap E) + \mu^*(A \setminus E) \text{ для всех } A \subset X \}
% \]
% является \( \sigma \)-алгеброй, а сужение \( \mu = \mu^*|_{\mathcal{M}} \) --- полной мерой на \( \mathcal{M} \). Доказывается, что все открытые множества лежат в \( \mathcal{M} \), следовательно, \( \mathcal{M} \) содержит борелевскую \( \sigma \)-алгебру.

% \textbf{3. Регулярность меры.}
% По построению мера \( \mu \) внешне регулярна:
% \[
% \mu(E) = \inf \{\mu(V) \colon V \text{ открыто}, E \subset V\} \quad \text{для всех борелевских } E.
% \]
% Внутренняя регулярность (для борелевских множеств конечной меры) устанавливается с использованием компактности:
% \[
% \mu(E) = \sup \{\mu(K) \colon K \text{ компакт}, K \subset E\}.
% \]

% \textbf{4. Представление функционала.}
% Пусть \( f \in C_c(X) \), \( f \geq 0 \). Для заданного \( \varepsilon > 0 \) выберем разбиение отрезка \( [0, \|f\|_\infty] \) точками \( 0 = t_0 < t_1 < \dots < t_n = \|f\|_\infty \) так, чтобы \( t_i - t_{i-1} < \varepsilon \). Определим множества
% \[
% E_i = \{ x \in X \colon t_{i-1} < f(x) \leq t_i \}, \quad i = 1, \dots, n.
% \]
% Каждое \( E_i \) есть разность двух открытых множеств, поэтому борелевское. Выберем компакты \( K_i \subset E_i \) и открытые \( V_i \supset E_i \) так, чтобы \( \mu(V_i \setminus K_i) < \varepsilon/n \). По лемме Урысона (используя локальную компактность и хаусдорфовость) существуют функции \( g_i \in C_c(X) \), \( 0 \leq g_i \leq 1 \), с носителем в \( V_i \), равные 1 на \( K_i \). Тогда
% \[
% \sum_{i=1}^n t_{i-1} g_i \leq f \leq \sum_{i=1}^n t_i g_i + \varepsilon \cdot 1_{K},
% \]
% где \( K = \operatorname{supp}(f) \). Применяя функционал \( \psi \) и используя его линейность и положительность, получаем оценки сверху и снизу для \( \psi(f) \). Переходя к пределу при \( \varepsilon \to 0 \), находим
% \[
% \psi(f) = \int_X f \, d\mu.
% \]
% Для произвольной \( f \in C_c(X) \) разложим её на положительную и отрицательную части.

% \textbf{5. Единственность меры.}
% Пусть \( \mu_1, \mu_2 \) --- регулярные борелевские меры, представляющие один и тот же функционал \( \psi \). Тогда для любого компакта \( K \) и любого открытого \( U \supset K \) существует функция \( f \in C_c(X) \) с \( 0 \leq f \leq 1 \), равная 1 на \( K \) и с носителем в \( U \). Отсюда
% \[
% \mu_1(K) \leq \int f \, d\mu_1 = \psi(f) = \int f \, d\mu_2 \leq \mu_2(U).
% \]
% Взяв инфимум по открытым \( U \supset K \), получим \( \mu_1(K) \leq \mu_2(K) \). Аналогично \( \mu_2(K) \leq \mu_1(K) \), так что \( \mu_1(K) = \mu_2(K) \). По внутренней регулярности меры совпадают на всех борелевских множествах.
% \end{proof}

% \begin{remark}
% Теорема Рисса-Маркова-Какутани является краеугольным камнем функционального анализа и теории меры. Она устанавливает взаимно однозначное соответствие между положительными линейными функционалами на пространстве непрерывных функций с компактным носителем и регулярными борелевскими мерами. В контексте предоставленных материалов эта теорема используется для представления линейных функционалов в доказательствах теорем аппроксимации (например, в теореме Цыбенко), что позволяет переходить от функционалов к интегральным представлениям.
% \end{remark}

% \subsection*{Теорема Хана-Банаха}
% \begin{theorem}[Хана-Банаха, аналитическая форма]
% Пусть \( X \) --- вещественное векторное пространство, \( p\colon X \to \mathbb{R} \) --- сублинейный функционал (т.е. \( p(x+y) \le p(x)+p(y) \) и \( p(tx)=t p(x) \) для \( t \ge 0 \)). Пусть \( M \subset X \) --- линейное подпространство, и \( f\colon M \to \mathbb{R} \) --- линейный функционал, удовлетворяющий условию \( f(x) \le p(x) \) для всех \( x \in M \). Тогда существует линейный функционал \( F\colon X \to \mathbb{R} \), продолжающий \( f \) (т.е. \( F|_M = f \)), такой что \( F(x) \le p(x) \) для всех \( x \in X \).
% \end{theorem}

% \subsection*{Преобразование Фурье}

% \begin{definition}[Преобразование Фурье в \( L^1(\mathbb{R}^n) \)]
% Для \( f \in L^1(\mathbb{R}^n) \) преобразование Фурье задается формулой:
% \[
% \hat{f}(\xi) = \int_{\mathbb{R}^n} f(x) e^{-2\pi i \xi \cdot x} \, dx.
% \]
% Тогда \( \hat{f} \in L^\infty \) и \( \|\hat{f}\|_\infty \leq \|f\|_1 \); более того, \( \hat{f} \in C_0(\mathbb{R}^n) \).
% \end{definition}

% \begin{definition}[Преобразование Фурье в \( L^2(\mathbb{R}^n) \)]
% Для \( f \in L^2(\mathbb{R}^n) \) преобразование Фурье определяется как предел в смысле \( L^2 \):
% \[
% \hat{f}(\xi) = \lim_{R \to \infty} \int_{|x| \leq R} f(x) e^{-2\pi i \xi \cdot x} \, dx.
% \]
% Оператор Фурье унитарен на \( L^2 \).
% \end{definition}

% Эти инструменты используются в гармоническом анализе и при доказательстве теорем об аппроксимации нейронными сетями.

\end{document}